%%% T08｜定积分及其应用 %%%


\section*{2020 年}


\textbf{选择题｜5 · 3 分}

设曲线 $y = x^2$ 与 $y = 4$ 所围成图形的面积为 $S$。则下列各式中错误的是( \quad )。\\
A. $S = 2 \int_0^2 (4 - x^2) dx$； \quad B. $S = 2 \int_0^4 \sqrt{y} dy$；\\
C. $S = 2 \int_0^2 x^2 dx$； \quad D. $S = 2 \int_0^4 \sqrt{x} dx$。


\textbf{计算题｜5 · 8 分}

设 $f(x) = \begin{cases} \frac{1}{1+x}, & x \ge 0 \\ \frac{1}{1+e^x}, & x < 0 \end{cases}$，求 $\int_0^2 f(x - 1) dx$。


\section*{2021 年}


\textbf{计算题｜5 · 8 分}

设 $F(x) = \int_0^{x^2} e^{-t^2} dt$，求（1）$F(x)$ 的拐点；（2）$\int_0^1 x^2 F'(x) dx$ 的值。


\textbf{应用题｜2 · 8 分}

在弹簧的拉伸过程中，拉力与弹簧的伸长量成正比。已知某弹簧拉伸 1 厘米（cm）需要的力是 3 牛顿（N），如果把该弹簧拉伸 4 厘米，计算需要做的功。


\textbf{证明题｜1 · 9 分}

已知 $f(x) = \begin{cases} x, & x < 0 \\ x+1, & x \ge 0 \end{cases}$，求 $F(x) = \int_{-1}^x f(t) dt \ (-1 \le x \le 1)$ 的表达式，并证明 $F(x)$ 在 $x = 0$ 处不可导。


\section*{2022 年}


\textbf{选择题｜5 · 3 分}

以下广义积分中发散的是（ \quad ）\\
A. $\int_e^{+\infty} \frac{1}{x\ln^2 x} dx$ \quad B. $\int_1^{+\infty} \frac{1}{x(1+x^2)} dx$ \quad C. $\int_{-1}^1 \frac{1}{\sin x} dx$ \quad D. $\int_{-1}^1 \frac{1}{\sqrt{1-x^2}} dx$


\textbf{选择题｜7 · 3 分}

矩形闸门宽 $a$ 米，高 $h$ 米，将其垂直放入水中，且闸门上沿与水面平齐，则闸门一侧所受的压力为（ \quad ）\\
A. $\rho g \int_0^h ax dx$ \quad B. $\rho g \int_0^a ax dx$ \quad C. $\rho g \int_0^h \frac{1}{2} ax dx$ \quad D. $\rho g \int_0^h 2ax dx$


\textbf{计算题｜7 · 8 分}

求曲线 $r = 1 - \cos\theta$ 的周长。


\textbf{应用题｜— · 9 分}

设 $0 < a < 2$，$D_1$ 是由抛物线 $y = 2x^2$ 和直线 $x = a, x = 2, y = 0$ 所围成的平面区域，$D_2$ 是由抛物线 $y = 2x^2$ 和直线 $x = a, y = 0$ 所围成的平面区域。\\
(1) 求 $D_1$ 绕 $x$ 轴旋转一周所形成的旋转体的体积 $V_1$ 和 $D_2$ 绕 $y$ 轴旋转一周所形成的旋转体的体积 $V_2$；\\
(2) 问 $a$ 为何值时，$V_1 + V_2$ 取得最大值？并求该最大值。


\textbf{证明题｜2 · 7 分}

设 $f(x)$ 是以 $\pi$ 为周期的连续函数，证明：$\int_0^{2\pi} (\sin x + x)f(x) dx = \int_0^\pi (2x + \pi)f(x) dx$。


\section*{2023 年}


\textbf{选择题｜5 · 3 分}

一容器的外表面由 $y = x^2 \ (0 \le x \le 1)$ 绕 $y$ 轴旋转而成，该容器盛满了水，若将容器内的水全部抽出，需要做的功为（ \quad ）。\\
A. $\rho g \int_0^1 \pi y^2 dy$ \quad B. $\rho g \int_0^1 \pi y(1-y^2) dy$\\
C. $\rho g \int_0^1 \pi y^2(1-y) dy$ \quad D. $\rho g \int_0^1 \pi y(1-y) dy$


\textbf{填空题｜3 · 3 分}

$\int_2^{+\infty} \frac{1}{x(\ln x)^2} dx = $ \underline{\hspace{3cm}}。


\textbf{计算题｜5 · 8 分}

设 $f(x) = \begin{cases} 1+x^2, & x < 0 \\ e^{-x}, & x \ge 0 \end{cases}$，求 $\int_{1/2}^2 f(x-1) dx$。


\textbf{计算题｜6 · 8 分}

设 $f(x) = \int_x^{x^2} e^{-t^2} dt$，计算 $\int_0^1 xf(x) dx$。


\textbf{计算题｜7 · 8 分}

求心形线 $r = a(1 + \cos\theta)$ 所围成的图形的面积 $S$。


\section*{2024 年}


\textbf{选择题｜4 · 3 分}

$\int_e^{+\infty} \frac{1}{x \ln^2 x} dx =$（ \quad ）\\
A. 0 \quad B. 1 \quad C. $\frac{1}{2}$ \quad D. 不存在


\textbf{计算题｜5 · 9 分}

计算定积分 $I = \int_0^1 2x \arctan x dx$。


\textbf{应用题｜2 · 9 分}

有一长方形闸门，形状与尺寸如下图所示，已知水面超过门顶 2m，求闸门所受水压力（注：水密度用 $\rho$ 表示，重力加速度用 $g$ 表示，取国际单位制单位）。


\section*{2025 年}


\textbf{选择题｜5 · 3 分}

积分 $\int_0^4 \frac{1}{x^2-x-6} dx$ （ \quad ）\\
A. $= 0$ \quad B. $= \frac{\pi}{4}$ \quad C. 发散 \quad D. $= \frac{1}{5}\ln\frac{3}{7}$


\textbf{应用题｜1 · 12 分}

（1）（6 分）求曲线 $y = 1 - \ln\cos x$ 上相应于 $0 \le x \le \frac{\pi}{6}$ 的一段弧的长度。\\
（2）（6 分）一游泳池宽 4 米、长 10 米、深 2 米，注满水后求游泳池四壁的压力（水密度为 $\rho$，重力加速度为 $g$，用国际单位制）。
